\section{Background and Setup}
\label{sec:bg}

\subsection{Workload}

We replay 200 coding-agent sessions drawn from the {\sc TraceLab} corpus
\cite{zhu2026tracelab}. Each session is a
sequence of rounds; a round presents a prompt of median 23{,}187 tokens
(p90 51{,}064; maximum 89{,}436, which alone is 88\% of the KV pool), and
between rounds the session waits on a tool before returning with a longer
context. Arrivals are Poisson at $\lambda=0.04\,\text{s}^{-1}$ over an
80-minute window, giving 999 steps in total. Two sessions of maximum size
cannot be co-resident: their combined contexts exceed the pool by 76\%.

The tool wait is the load's defining property and the one most easily
misread. Its \emph{median} is 0.35\,s, which invites the conclusion that
sessions never leave the engine. Its \emph{mean} is 16.85\,s, 22\% of waits
exceed ten seconds, and the maximum is 300\,s. We verified that our replay
harness honors the trace faithfully: across 799 round transitions the measured
gap is 16.85\,s mean against a demanded 15.82\,s, a ratio of 1.001. The
corpus's own characterization makes the same point qualitatively---it reports
``diverse and heavily-tailed tool calls'' and calls these gaps
``human-paced''---but gives no numeric mean \cite{zhu2026tracelab}, so the
16.85\,s above is our measurement over the released trace, not a figure we
quote from that paper.

\subsection{Engine and the binding constraint}

We use SGLang v0.5.10 \cite{zheng2024sglang} with Qwen2.5-Coder-7B-Instruct
\cite{qwen2coder} extended to 96K context,
on one RTX 3090 (24\,GB). KV is paged \cite{kwon2023pagedattention}; at \texttt{mem-fraction-static} 0.85
the pool is 101{,}432 tokens. Hierarchical caching (HiCache) is enabled throughout, with a
host tier of twice the device pool. Every run is guarded: if the pool does not
come out at $101{,}432 \pm 20$---which happens when another tenant holds
memory on the card---the run is aborted rather than recorded, because SGLang
sizes the pool from free memory at startup and a shrunken pool silently
produces incomparable numbers.

The relevant consequence is that the pool holds only about four median
sessions. Concurrency is therefore capacity-bound, not scheduler-bound, and
this is the premise on which everything below rests.

\subsection{The controller under test}

\sys{} is an admission controller that reads engine telemetry and issues
\emph{named} permits rather than a scalar cap: a session is admitted by
appearing in a permit file, and the runner refuses to start a round for a
session that is not listed. This design choice matters for the comparison,
because it is what allows per-session actions at all---including the pause
primitive we exploit in \S\ref{sec:results:gate}. It also means that a session,
once admitted, holds its permit for its entire lifetime and its KV is
preserved across tool waits. We return to this point in \S\ref{sec:discussion},
because we believe it is why published gains from state-preserving suspension
do not reappear here.

Each control cycle (2\,s) the controller computes
\begin{equation}
  \Pi \;=\; \max(\rho, \textstyle\sum_i c_i) \;+\; \textstyle\sum_i g_i ,
  \label{eq:proj}
\end{equation}
where $\rho$ is the engine's reported occupancy, $c_i$ the context length of
live session $i$, and $g_i$ its forecast growth over the next
\texttt{horizon} rounds. A session is admitted when $\Pi + n \le B$, with $n$
the candidate's estimated need and $B$ the budget. The three terms are
respectively the measured physical occupancy, a reservation for sessions
currently parked in a tool wait (whose KV is evictable and therefore absent
from $\rho$), and a reservation for future growth.

\subsection{Baselines and the verdict criterion}

Baselines are static caps: at most $N$ sessions resident, $N \in \{1,\dots,5\}$,
with and without HiCache. We compare on two axes: end-to-end wall clock
(time from the first submission to the last completion) and SLO attainment
(fraction of requests whose time-to-first-token is under 10\,s).

The verdict criterion deserves care, because ``faster than cap-$N$'' is not
the question. The question is whether a policy reaches a point \emph{above}
the static front. A point that lies merely between two static caps is not a
win: mixing cap-$N$ and cap-$M$ by time slice approximates any point on the
line between them. We flag this segment argument as an argument, not a
measurement---we did not run the time-sliced policies, and carrying queue state
and cache contents across a switch could bend the line. It is the criterion we
apply, and it is conservative in the direction that matters: it can only make
our negative result harder to claim, never easier. So we take the upper-left
envelope of the static points and report each controller's signed distance
above or below it. Every verdict in
this paper is produced by one script from the raw logs, and the script's own
first version was wrong---it extrapolated past the end of the envelope and
declared the slowest static cap a winner. We fixed it and report the episode
because it is the failure mode this criterion is most prone to.
